\chapter{Michael Atiyah \& Isadore Singer (2004)}

\section{Biographical Sketch}



\subsection{Sir Michael Francis Atiyah}

Michael Atiyah was born on 22 April 1929 in London, England. He grew up in Sudan and Lebanon, returning to England for his education. He attended the University of Manchester and later Trinity College, Cambridge, where he completed his PhD under the supervision of W. V. D. Hodge in 1955. He held positions at the Institute for Advanced Study (Princeton), Oxford University (Savilian Professor of Geometry), and the University of Cambridge (Master of Trinity College). He served as President of the Royal Society (1990--1995) and was knighted in 1983. Atiyah received the Fields Medal in 1966 and the Abel Prize (jointly with Singer) in 2004.

\subsection{Isadore Manuel Singer}

Isadore Singer was born on 3 May 1924 in Detroit, Michigan, USA. He attended the University of Michigan for his undergraduate studies and the University of Chicago for his PhD, which he completed in 1950 under the supervision of Irving Segal. He held positions at the University of California, Berkeley, the Massachusetts Institute of Technology, and the University of California, Los Angeles. Singer received the Abel Prize (jointly with Atiyah) in 2004 and the National Medal of Science in 1983.

\section{High-Level View for a General Audience}

\subsection{What Did Atiyah and Singer Do?}

Imagine you have a beach ball with a gentle breeze blowing across its surface. At every point on the ball, the wind has a direction and a speed---this is a \emph{vector field}. The hairy ball theorem, a classic result in topology, says that you cannot comb a hairy ball flat without creating a cowlick: every continuous vector field on a sphere must vanish at at least one point. In other words, somewhere on Earth, the wind is perfectly still.

This is a simple example of a much deeper relationship between the \emph{local} behavior of a system (how the wind blows near each point) and the \emph{global} structure of the space it lives on (the fact that the sphere has no holes of a certain type). The Atiyah--Singer index theorem is a breathtakingly general version of this idea. It says that for a wide class of differential equations---equations that describe how things change---there is a deep equality between:

\begin{enumerate}
\item The number of independent solutions to the equation (an analytic, local quantity), and
\item A topological invariant of the space on which the equation is defined (a geometric, global quantity).
\end{enumerate}

\subsection{Why Does It Matter?}

The Atiyah--Singer index theorem is often described as one of the deepest results of twentieth-century mathematics. Its power comes from its universality: it applies to elliptic differential operators on any compact manifold. This means it has applications to:

\begin{itemize}
\item \textbf{Physics}: The theorem is essential for understanding anomalies in quantum field theory, the geometry of string theory, and the mathematical structure of gauge theories.

\item \textbf{Topology}: The theorem provides a way to compute topological invariants by solving differential equations, and conversely, to compute analytic invariants from topology.

\item \textbf{Number Theory}: The theorem has connections to the arithmetic of quadratic forms, modular forms, and the Langlands program.
\end{itemize}

In simple terms: the Atiyah--Singer index theorem tells us that the solutions to differential equations are profoundly constrained by the shape of the space they live on. It is a bridge between analysis and topology that has shaped modern mathematics.

\section{The Origin of the Problem}

\subsection{The Index of an Operator}

The story begins with a simple linear algebra fact. If $T: V \to W$ is a linear map between finite-dimensional vector spaces, then:
\[
\dim \ker(T) - \dim \operatorname{coker}(T) = \dim V - \dim W.
\]
The left-hand side is called the \emph{index} of $T$. It measures the extent to which $T$ is invertible. The right-hand side depends only on the dimensions of the domain and codomain.

For differential operators on manifolds, the situation is far more subtle. If $D: C^\infty(E) \to C^\infty(F)$ is a linear differential operator between spaces of smooth sections of vector bundles $E$ and $F$, then the kernel and cokernel of $D$ may be infinite-dimensional. However, for \emph{elliptic} operators on a compact manifold, both the kernel and cokernel are finite-dimensional, and the index:
\[
\operatorname{ind}(D) = \dim \ker(D) - \dim \operatorname{coker}(D)
\]
is a well-defined integer.

\subsection{Elliptic Operators and Their Significance}

An elliptic differential operator is one that satisfies a certain positivity condition on its symbol. The prototypical example is the Laplacian:
\[
\Delta = -\frac{\partial^2}{\partial x_1^2} - \frac{\partial^2}{\partial x_2^2} - \cdots - \frac{\partial^2}{\partial x_n^2}.
\]

Elliptic operators arise naturally in many contexts:
\begin{itemize}
\item The Laplace equation $\Delta u = 0$ describes steady-state heat distribution.
\item The Cauchy--Riemann equations $\partial f / \partial \bar z = 0$ define holomorphic functions.
\item The Dirac equation $D\!\!\!\!/ \ \psi = 0$ describes relativistic electrons.
\item The Hodge Laplacian $\Delta = d\delta + \delta d$ relates to cohomology\index{Cohomology}.
\end{itemize}

A fundamental property of elliptic operators on compact manifolds is the \emph{elliptic estimate}: if $Du = f$ and $u$ is bounded in a suitable norm, then $u$ is smooth. This leads to the Fredholm property: elliptic operators on compact manifolds have finite-dimensional kernel and cokernel, and closed range.

\subsection{Early Index Theorems}

Before Atiyah and Singer, several special cases of the index theorem were known:

\begin{itemize}
\item \textbf{The Chern--Gauss--Bonnet Theorem}: For the de Rham complex on a compact oriented manifold, the index of the Hodge--Dirac operator $d + d^*$ is the Euler characteristic $\chi(M)$. The Chern--Gauss--Bonnet theorem expresses $\chi(M)$ as the integral of a curvature form:
\[
\chi(M) = \int_M \frac{1}{(2\pi)^{n/2}} \operatorname{Pf}(\Omega),
\]
where $\Omega$ is the curvature form of any Riemannian metric.

\item \textbf{The Hirzebruch Signature Theorem}: For the signature complex on a compact oriented $4k$-dimensional manifold, the index is the signature $\sigma(M)$. Hirzebruch (1954) showed that:
\[
\sigma(M) = \int_M L(p_1, \ldots, p_k),
\]
where $L$ is the Hirzebruch $L$-polynomial in the Pontryagin classes.

\item \textbf{The Riemann--Roch Theorem}: For the $\bar\partial$ complex on a Riemann surface, the index is $1 - g$, where $g$ is the genus. The classical Riemann--Roch theorem was generalized by Hirzebruch to higher dimensions using the Todd genus.
\end{itemize}

These results were striking: in each case, an analytic invariant (the index of a differential operator) equals a topological invariant (integral of a characteristic class). Atiyah and Singer set out to understand this phenomenon in complete generality.

\section{Deep Insight into the Problem}

\subsection{The K-Theoretic Formulation}

Atiyah and Singer's deep insight was that the index of an elliptic operator should be understood in terms of $K$-theory, a generalized cohomology\index{Cohomology} theory that Atiyah and Hirzebruch had developed in the early 1960s.

The symbol $\sigma(D)$ of an elliptic operator $D$ defines a class in $K(T^*M)$, the $K$-theory of the cotangent bundle. Atiyah and Singer showed that the index of $D$ depends only on this symbol class and that the map from $K(T^*M)$ to $\mathbb{Z}$ given by the index is a homomorphism.

This observation has two crucial consequences:
\begin{enumerate}
\item The index is a topological invariant---it depends only on the symbol class, not on the specific operator.
\item The index factors through the $K$-theory of $T^*M$, allowing the use of topological methods to compute it.
\end{enumerate}

\subsection{The Index Theorem}

The Atiyah--Singer index theorem states:
\[
\operatorname{ind}(D) = (-1)^{n(n+1)/2} \int_{T^*M} \operatorname{ch}(\sigma(D)) \cdot \pi^* \operatorname{Td}(TM \otimes \mathbb{C}),
\]
where $\operatorname{ch}$ is the Chern character, $\operatorname{Td}$ is the Todd class, and $\pi: T^*M \to M$ is the projection.

Equivalently, for an elliptic differential operator $D$ on a compact $n$-dimensional manifold $M$:
\[
\operatorname{ind}(D) = \int_M \operatorname{ch}(\sigma(D)) \cdot \widehat A(M) \quad \text{(for Dirac-type operators)},
\]
or more generally:
\[
\operatorname{ind}(D) = \int_M \operatorname{ch}(\sigma(D)) \cdot \operatorname{Td}(M).
\]

\subsection{The Proof Strategy}

The proof of the index theorem proceeds in several steps:

\begin{enumerate}
\item \textbf{K-Theoretic Formulation}: Show that the index depends only on the $K$-theory class of the symbol, giving a homomorphism $\operatorname{ind}: K(T^*M) \to \mathbb{Z}$.

\item \textbf{Axiomatics}: Show that the index is uniquely determined by a small set of axioms: normalization, excision, and homotopy\index{Homotopy groups} invariance. Both the analytic index (the left-hand side) and the topological index (the right-hand side) satisfy these axioms.

\item \textbf{Equality}: Since both sides satisfy the same axioms, they must be equal. The proof uses a clever embedding argument to reduce the problem to the case of a sphere, where explicit computation verifies the equality.

\item \textbf{Embedding and Tubular Neighborhood}: Embed $M$ into $\mathbb{R}^N$ for large $N$, extend the operator to a neighborhood, and use the product structure of $\mathbb{R}^N$ to reduce the computation to a standard case.
\end{enumerate}

\section{Biggest Contributions and New Tools}

\subsection{The Index Theorem Itself}

The Atiyah--Singer index theorem is the crowning achievement of the collaboration. It unified and vastly generalized the known index theorems (Gauss--Bonnet, Hirzebruch signature, Riemann--Roch) and provided a single conceptual framework for understanding them all.

\subsection{$K$-Theory}

Atiyah and Hirzebruch developed $K$-theory in the late 1950s and early 1960s as a tool for studying vector bundles. $K$-theory is a generalized cohomology\index{Cohomology} theory defined by:
\[
K^0(X) = \text{ Grothendieck group of vector bundles over } X.
\]

$K$-theory had already been used to great effect in:
\begin{itemize}
\item The solution of the Hopf invariant one problem (Adams, 1960).
\item The computation of the number of linearly independent vector fields on spheres (Adams, 1962).
\item The Bott periodicity theorem (Bott, 1959), which states that $K(S^{n+2}) \cong K(S^n)$.
\end{itemize}

Atiyah and Singer made $K$-theory an essential part of the index theorem, using it to encode the symbol of an elliptic operator in purely topological terms.

\subsection{The Elliptic Genus}

The Atiyah--Singer index theorem led to the discovery of the \emph{elliptic genus}, a genus that associates to every compact oriented manifold a modular form. The elliptic genus, discovered by Ochanine, Landweber, and Stong in the 1980s, is the index of a certain Dirac operator on the loop space.

The elliptic genus has deep connections to:
\begin{itemize}
\item Modular forms and number theory.
\item String theory, where it appears as the partition function of the superstring.
\item Mirror symmetry in algebraic geometry.
\end{itemize}

\subsection{The $\widehat A$-Genus and the Dirac Operator}

The $\widehat A$-genus, defined by:
\[
\widehat A(M) = \prod_{i=1}^{n/2} \frac{x_i/2}{\sinh(x_i/2)},
\]
where $x_i$ are the Pontryagin roots, plays a central role in the index theorem for the Dirac operator.

Atiyah and Singer showed that the index of the Dirac operator on a compact spin manifold is the $\widehat A$-genus:
\[
\operatorname{ind}(D\!\!\!\!/) = \widehat A(M).
\]

This result has profound consequences:
\begin{itemize}
\item It provides a topological obstruction to the existence of positive scalar curvature metrics (Lichnerowicz, Hitchin).
\item It connects to the theory of spinors and the geometry of manifolds with special holonomy.
\item It plays a central role in the analysis of anomalies in quantum field theory.
\end{itemize}

\subsection{Equivariant Index Theory}

Atiyah and Singer developed an equivariant version of the index theorem for manifolds with a group action. The equivariant index theorem states that for a compact Lie group\index{Lie groups} $G$ acting on a manifold $M$ and commuting with an elliptic operator $D$, the index takes values in the representation ring $R(G)$ and can be computed by a fixed-point formula.

The equivariant index theorem has applications to:
\begin{itemize}
\item The Weyl character formula for representations of compact Lie groups\index{Lie groups}.
\item The geometry of flag varieties and homogeneous spaces.
\item The computation of cohomology of line bundles on projective varieties.
\end{itemize}

\subsection{The Lefschetz Fixed-Point Formula}

The Atiyah--Singer equivariant index theorem implies a powerful generalization of the Lefschetz fixed-point formula. For a map $f: M \to M$ that preserves an elliptic operator, the Lefschetz number:
\[
L(f) = \sum_i (-1)^i \operatorname{Tr}(f^*|_{H^i(M)})
\]
can be expressed as a sum over fixed points of $f$ of certain local contributions.

The Holomorphic Lefschetz formula, proved by Atiyah and Bott using these methods, is a far-reaching generalization that applies to holomorphic maps on complex manifolds.

\section{Impact of the Discovery in Science and Technology}

\subsection{Impact on Mathematics}

The Atiyah--Singer index theorem is one of the most far-reaching results in modern mathematics:

\begin{itemize}
\item \textbf{Global Analysis}: The theorem established the centrality of elliptic operators and their indices in the study of manifolds. It showed that analytic methods (solving differential equations) can reveal topological information, and vice versa.

\item \textbf{Algebraic Geometry}: The Riemann--Roch theorem for projective varieties (Hirzebruch--Riemann--Roch) and its generalization by Grothendieck are consequences of the index theorem. The Grothendieck--Riemann--Roch theorem is a key tool in intersection theory and motivic cohomology.

\item \textbf{Representation Theory}: The equivariant index theorem provides a geometric approach to representation theory. The Atiyah--Bott fixed-point formula is used to compute character formulas and cohomology of homogeneous spaces.

\item \textbf{Differential Geometry}: The index theorem for the Dirac operator provides obstructions to the existence of metrics with positive scalar curvature. It also has applications to the geometry of spin manifolds and manifolds with special holonomy.

\item \textbf{Topology}: The index theorem provides a way to compute topological invariants using analysis. For example, the signature of a manifold can be computed as the index of the signature operator, which can be evaluated using characteristic classes.
\end{itemize}

\subsection{Impact on Theoretical Physics}

The Atiyah--Singer index theorem has had a transformative impact on theoretical physics:

\begin{itemize}
\item \textbf{Gauge Theory}: The index theorem is used to compute the number of zero modes of the Dirac operator in gauge fields, which is essential for understanding instantons and monopoles. The Atiyah--Singer index theorem for the Dirac operator in the presence of a gauge field gives the Atiyah--Patodi--Singer index theorem, which includes boundary contributions.

\item \textbf{Quantum Anomalies}: The cancellation of chiral anomalies in quantum field theory---essential for the consistency of the Standard Model---is understood using the index theorem. The 't Hooft anomaly matching conditions are a consequence of index-theoretic considerations.

\item \textbf{String Theory}: The elliptic genus, derived from the index theorem, is the partition function of the superstring. The index theorem on loop spaces, developed by Atiyah, Witten, and others, is central to string theory.

\item \textbf{Topological Quantum Field Theory}: The Atiyah--Singer index theorem inspired Atiyah's axiomatization of topological quantum field theory (TQFT). The Jones polynomial and other knot invariants can be understood using TQFT and the index theorem.

\item \textbf{Seiberg--Witten Theory}: The Seiberg--Witten equations (1994), which revolutionized four-dimensional topology, involve the Dirac operator coupled to a gauge field. The index theorem is essential for understanding the moduli space of solutions.
\end{itemize}

\subsection{Impact on Technology}

While the index theorem is a highly abstract mathematical result, it has influenced technological developments:

\begin{itemize}
\item \textbf{Signal Processing}: The theory of pseudodifferential operators, which is essential for the index theorem, has applications to signal processing and image analysis. The symbolic calculus of pseudodifferential operators is used in the analysis of partial differential equations arising in engineering.

\item \textbf{Computer Vision}: The theory of the heat kernel, which is used in the heat kernel proof of the index theorem, has applications to shape recognition and computer vision. The heat kernel signature is used to analyze the geometry of shapes.

\item \textbf{Cryptography}: The arithmetic of elliptic curves\index{Elliptic curves}, which is related to the index theorem through the theory of modular forms, is the basis of elliptic curve\index{Elliptic curves} cryptography.

\item \textbf{Machine Learning}: The geometric methods inspired by the index theorem, including the study of manifolds and their invariants, have found applications in machine learning, particularly in manifold learning and dimensionality reduction.
\end{itemize}

\section{Current Developments}

\subsection{The Atiyah--Patodi--Singer Index Theorem}

The extension of the index theorem to manifolds with boundary, developed by Atiyah, Patodi, and Singer in the 1970s, remains an active area of research:

\begin{itemize}
\item \textbf{Boundary Value Problems}: The APS index theorem provides a framework for studying elliptic boundary value problems and their indices. The eta invariant, which measures the spectral asymmetry of the boundary operator, has deep connections to topology and number theory.

\item \textbf{Applications to Geometry}: The APS index theorem is used to study the geometry of manifolds with boundary, including the topology of four-manifolds and the geometry of hyperbolic manifolds.
\end{itemize}

\subsection{The Index Theorem in Noncommutative Geometry}

Connes's noncommutative geometry extends the index theorem to noncommutative spaces:

\begin{itemize}
\item \textbf{The Connes--Moscovici Index Theorem}: This generalizes the Atiyah--Singer index theorem to foliations and noncommutative spaces. It is used to study the geometry of quotient spaces and the arithmetic of number fields.

\item \textbf{The Local Index Formula}: Connes and Moscovici proved a local index formula in noncommutative geometry, expressing the index of a spectral triple in terms of local expressions involving the heat kernel.
\end{itemize}

\subsection{The Index Theorem and the Langlands Program}

The geometric Langlands program, which relates the representation theory of reductive groups to the geometry of moduli spaces, uses the index theorem in essential ways:

\begin{itemize}
\item \textbf{The Geometric Satake Correspondence}: The Satake isomorphism, which relates spherical functions to representations of the Langlands dual group, has a geometric proof using the index theorem.

\item \textbf{Hecke Eigensheaves}: The construction of Hecke eigensheaves, which are the geometric analogs of automorphic forms, uses the index theorem for the $\bar\partial$ operator on the moduli space of bundles.
\end{itemize}

\subsection{The Index Theorem and Mirror Symmetry}

Mirror symmetry, one of the most active areas of current research in mathematical physics, has deep connections to the index theorem:

\begin{itemize}
\item \textbf{The Elliptic Genus and Mirror Symmetry}: The elliptic genus, which is the index of the Dirac operator on loop space, is a key tool in the study of mirror symmetry. The genus is invariant under mirror symmetry, providing a test of mirror relationships.

\item \textbf{Gromov--Witten Invariants}: The Gromov--Witten invariants, which count holomorphic curves in algebraic varieties, are related to the index of the $\bar\partial$ operator on the moduli space of maps.
\end{itemize}

\subsection{Computational Aspects}

The computation of indices using the Atiyah--Singer theorem is an active area of computational mathematics:

\begin{itemize}
\item \textbf{Algorithms for Characteristic Classes}: The computation of Pontryagin, Chern, and Stiefel--Whitney classes, which appear in the index formula, is implemented in software packages like Macaulay2 and SageMath.

\item \textbf{The Index in Mathematical Physics}: Software for computing indices of Dirac operators in gauge theory, including instanton and monopole moduli spaces, is used in mathematical physics research.

\item \textbf{Machine Learning}: Neural networks have been used to learn topological invariants from geometric data, providing a computational approach to the index theorem.
\end{itemize}

\section{Detailed Analysis of Key Results}

\subsection{The Hirzebruch--Riemann--Roch Theorem}

The Hirzebruch--Riemann--Roch theorem, which preceded and inspired the index theorem, states that for a holomorphic vector bundle $V$ on a compact complex manifold $M$:
\[
\chi(M, V) = \int_M \operatorname{ch}(V) \cdot \operatorname{Td}(M),
\]
where $\chi(M, V) = \sum_i (-1)^i \dim H^i(M, \mathcal{O}(V))$ is the Euler characteristic of $V$.

Atiyah and Singer showed that this is a special case of their theorem: the index of the $\bar\partial$ operator on $M$ with coefficients in $V$ is exactly $\chi(M, V)$.

\subsection{The Signature Theorem}

For a compact oriented $4k$-dimensional manifold $M$, the signature $\sigma(M)$ is the index of the signature operator, and the index theorem gives:
\[
\sigma(M) = \int_M L(p_1, \ldots, p_k),
\]
where $L$ is the Hirzebruch $L$-polynomial.

For $4$-manifolds, this becomes:
\[
\sigma(M) = \frac{1}{3} \int_M p_1(M),
\]
which is the Hirzebruch signature theorem. This is a key result in the classification of four-manifolds.

\subsection{The $\widehat A$-Genus and the Dirac Operator}

For a compact spin manifold $M$ with Dirac operator $D\!\!\!\!/$, the index theorem gives:
\[
\operatorname{ind}(D\!\!\!\!/) = \widehat A(M).
\]

A striking consequence is that if $M$ admits a Riemannian metric with positive scalar curvature, then the $\widehat A$-genus must vanish (the Lichnerowicz theorem). This provides a topological obstruction to positive scalar curvature metrics.

\section{Conclusion}

The Atiyah--Singer index theorem stands as one of the greatest achievements of twentieth-century mathematics. It unified diverse areas of mathematics---analysis, topology, geometry, and algebra---under a single conceptual framework. The theorem showed that the number of solutions to an elliptic differential equation on a manifold is determined by the topology of the manifold itself, a deep and beautiful result.

The impact of the theorem extends far beyond mathematics. It has shaped modern theoretical physics, providing essential tools for understanding gauge theories, anomalies, and string theory. It continues to inspire new developments in noncommutative geometry, the Langlands program, and mirror symmetry.

Atiyah and Singer were awarded the Abel Prize in 2004 for their discovery and proof of the index theorem, and for its profound impact on mathematics. As the citation notes, the theorem "brought together topology, geometry, and analysis for the benefit of all three."


\section{Behind the Breakthrough: The Human Story}
Atiyah and Singer met at the Institute for Advanced Study in Princeton. They realized they were attacking the exact same problem from opposite ends: Atiyah from algebraic geometry/topology and Singer from analysis/differential operators. Their collaboration lasted for decades and completely unified the two fields.

\section{Pedagogical Metaphors and Lineage}
\subsection{Signature Metaphor}
``A mathematical scale that weighs a geometric shape to tell you how many solutions a differential equation has, without you ever having to solve the equation.''

\subsection{Mathematical Lineage}
Atiyah was advised by William Hodge (originator of Hodge theory\index{Hodge theory}). Singer was advised by Richard Duffin. Atiyah went on to advise Simon Donaldson and Nigel Hitchin.

\section{Applied Science and Computational Algorithms}
\subsection{Key Takeaways for Applied Science}
The index theorem is widely used in gauge field theories and quantum anomalies in theoretical physics, enabling physicists to compute vacuum states, charge conservation violations, and anomalies in string theory and quantum gravity.

\subsection{Algorithms and Numerical Implementations}
Lattice QCD (Quantum Chromodynamics) simulations utilize discretized Dirac operator schemes (such as overlap fermions) to numerically calculate the index and determine topological properties of the vacuum on a discrete space-time grid.

\section{The Research Frontier}
Extending index theory to non-commutative spaces (Alain Connes' non-commutative geometry) and infinite-dimensional manifolds (loop spaces and Dirac operators in quantum field theory) represents the modern research frontier.

\section{Guided Exercises and Challenges}
\begin{enumerate}[label=\textbf{\arabic*.}]
  \item State the classical Gauss-Bonnet theorem and show how it expresses the Euler characteristic as the index of a de Rham-Dirac operator.
  \item Describe the difference between the analytical index and the topological index of an elliptic differential operator.
  \item Explain how the Atiyah-Patodi-Singer (APS) index theorem handles manifolds with boundary using the eta invariant.
\end{enumerate}

\section{Curriculum Map and Further Reading}
For students wishing to delve deeper into the mathematical machinery underlying these breakthroughs, the following learning path is recommended:
\begin{itemize}
  \item H. Blaine Lawson and Marie-Louise Michelsohn, \emph{Spin Geometry}, Princeton University Press.
  \item N. Berline, E. Getzler, and M. Vergne, \emph{Heat Kernels and Dirac Operators}, Springer.
\end{itemize}

\section*{Bibliography}

\begin{enumerate}
\item M. F. Atiyah and I. M. Singer, "The Index of Elliptic Operators on Compact Manifolds," Bulletin of the American Mathematical Society 69 (1963), 422--433.
\item M. F. Atiyah and I. M. Singer, "The Index of Elliptic Operators I," Annals of Mathematics 87 (1968), 484--530.
\item M. F. Atiyah and I. M. Singer, "The Index of Elliptic Operators II," Annals of Mathematics 87 (1968), 531--545.
\item M. F. Atiyah and I. M. Singer, "The Index of Elliptic Operators III," Annals of Mathematics 87 (1968), 546--604.
\item M. F. Atiyah and I. M. Singer, "The Index of Elliptic Operators IV," Annals of Mathematics 92 (1970), 119--138.
\item M. F. Atiyah and I. M. Singer, "The Index of Elliptic Operators V," Annals of Mathematics 93 (1971), 139--149.
\item M. F. Atiyah, R. Bott, and V. K. Patodi, "On the Heat Equation and the Index Theorem," Inventiones Mathematicae 19 (1973), 279--330.
\item M. F. Atiyah, V. K. Patodi, and I. M. Singer, "Spectral Asymmetry and Riemannian Geometry I--III," Mathematical Proceedings of the Cambridge Philosophical Society 77--79 (1975--1976).
\item P. B. Gilkey, "Invariance Theory, the Heat Equation, and the Atiyah--Singer Index Theorem," CRC Press, 1995.
\item N. Berline, E. Getzler, and M. Vergne, "Heat Kernels and Dirac Operators," Springer, 1992.
\end{enumerate}
